%Because tokens are longest possible valid string and because we first match keywords, then identifiers, true(); is a parse error.
%ikke derfor men ja det er en parse error.

%TODO: Husk at nævn afvigelser fra tests.


% Idéer til afsnit
% 1. Overordnet om parser
%   - performs syntactic analysis of VVPL programs
%   - Uses tokens from scanner to construct an AST
%   - What this AST consists of (Expr, Stmt, polymorphism)
%   - A tree that matches the grammatical structure of the language
%   - A recursive descent parser. 
%   - Noget med at forklare den overordnede rekursive natur og hvordan parser blot mapper tokens til den rekursive grammar
% 2. Error recovery (synchronization)
%   - Tænker dette er noget af det lidt fede ved parseren og god blanding af teori og det tekniske
%   - Kan genbruge mange af idéerne fra monitoring and logging afsnttet, som ikke længere skal bruges.
% 3. Detaljer
%   - Flet nogle af detaljerne ind undervejs. 
%   - Vigtigste er: håndtering af casts, consume rapporterer previous og måske true()



The parser performs syntactic analysis of VVPL programs.

Recursive-descent parser based directly on the VVPL grammar.

Consumes a list of Tokens produced by the Scanner and constructs an abstract syntax tree (AST).
This conists of Expr and Stmt nodes.
We use polymorphism to allow the use of the more detailed subclasses.


Detaljer som måske er værd at nævne:
- Wrapper block stmt. Kan ikke helt huske hvorfor dette var smart, men mener også bogen gør det.
- Håndtering af casts. Vi har valgt at have en seperat CAST node.
- Vi håndterer true() som en parse error selvom det er beskrevet som en scope error i projektbeskrivelsen
- Vi bruger en custom "Param" klasse til at håndterer funktioners parametre.
- Consume rapporterer error på previous().
    Er fint fordi vi altid matcher inden vi kalder consume. Så vi ved altid at mindst én token er
    fundet inden vi kalder consume. Så vi kan ikke komme out of bounds.
    Valget er taget for at manglende semikoloner bliver rapporteret på den rigtige linje.
    Lidt ulogisk i forhold til at vi ikke rapporterer error på den token der faktisk skabte fejlen.
    Men pga. vi kun rapporterer linje, så virker det godt i dette tilfælde.